\chapter{Language models, from first principles}
\label{ch:transformers}

This chapter builds up exactly enough machinery to understand what SilkomeGPT
is. If you already know what a decoder-only transformer is, skip to
Section~\ref{sec:decoding}, which covers the decoding settings --- those turn
out to matter more for this project than the architecture does.

\section{The one thing a language model does}

A language model does one thing: given a piece of text, it assigns a
probability to every possible next chunk of text.

That is genuinely all. Everything else --- chat, translation, protein design ---
is built by choosing what to put in and how to read what comes out.

Write the input as a sequence of tokens $t_1, t_2, \dots, t_n$. The model
computes
\[
  P(t_{n+1} \mid t_1, \dots, t_n)
\]
a probability for each of the $\approx$50{,}000 possible next tokens. To
generate text you sample one from that distribution, append it, and repeat.
This is called \term{autoregressive} generation: each new token is conditioned
on everything already written, including what the model itself just wrote.

\begin{keybox}{Why this is enough to design proteins}
If you train the model on text of the form
\begin{center}
\code{[a description of what I want][the protein that satisfies it]}
\end{center}
then at generation time you write the description, and the machinery that
predicts ``the most likely continuation'' has no choice but to produce a
protein that fits it. Design becomes autocomplete. The trick is entirely in the
format of the training text.
\end{keybox}

\section{Tokens: how text becomes numbers}

Neural networks consume numbers, not letters. A \term{tokenizer} chops text
into pieces and assigns each piece an integer ID. The \term{vocabulary} is the
set of pieces it knows.

The naive choice for proteins is one token per amino acid, giving a vocabulary
of twenty. SilkomeGPT does something more efficient, called \term{byte-pair
encoding} (BPE), and the reason matters for silk in particular.

\begin{keybox}{Byte-pair encoding, by example}
BPE learns its vocabulary from data by repeatedly merging the most frequent
adjacent pair. Suppose the training data is full of the fragment
\aaseq{GAGAGA}. Then:

\begin{enumerate}[leftmargin=*,itemsep=1pt]
  \item Start with single letters: \aaseq{G}, \aaseq{A}.
  \item \aaseq{G}+\aaseq{A} is the most frequent adjacent pair $\rightarrow$
        merge into a new token \aaseq{GA}.
  \item Now \aaseq{GA}+\aaseq{GA} is frequent $\rightarrow$ merge into
        \aaseq{GAGA}.
  \item Continue until the vocabulary reaches the target size.
\end{enumerate}

The result: common motifs become \emph{single tokens}. A poly-alanine run
\aaseq{AAAAAAAA} might be one token rather than eight.
\end{keybox}

This is why the numbers in this project look odd at first glance. The paper's
worked-example sequence is 635 amino acids long but only 198 tokens. That is
about 3.2 residues per token, and it happens because spidroins are extremely
repetitive --- exactly the situation BPE compresses well. SilkomeGPT's tokenizer
was trained on $\approx$800{,}000 protein sequences with a vocabulary of up to
50{,}000.

\begin{warnbox}{A trap this creates}
Because compression depends on how repetitive a sequence is, token count is not
proportional to residue count. A highly repetitive 2000-residue spidroin may
tokenise to fewer tokens than a non-repetitive 800-residue one. Any reasoning
about ``context limits'' has to be done in tokens, not residues.
\end{warnbox}

\section{Attention, and what makes a transformer}

The architecture underneath is a \term{transformer}. Its central mechanism is
\term{self-attention}, and the idea is simpler than its reputation.

When the model processes a token, it needs context from other tokens. Attention
lets every position look at every other position and decide how much each one
matters, for this particular purpose, right now.

\begin{figure}[htbp]
\centering
\begin{tikzpicture}[font=\small,
  tok/.style={draw=greyc!70,fill=boxbg,rounded corners=2pt,
              minimum width=1.25cm,minimum height=0.7cm},
  cur/.style={tok,draw=ourscol,fill=ourscol!12,line width=1pt},
]
\node[tok] (t1) at (0,0)   {\aaseq{GGA}};
\node[tok] (t2) at (1.5,0) {\aaseq{GQGG}};
\node[tok] (t3) at (3.0,0) {\aaseq{AAAAAA}};
\node[tok] (t4) at (4.5,0) {\aaseq{GYG}};
\node[cur] (t5) at (6.0,0) {\aaseq{QGG}};

\node[font=\scriptsize,below=2pt of t5,text=ourscol]
  {current position};

% attention arcs, thickness = weight
\draw[-{Stealth[length=2mm]},ourscol,line width=2.2pt,opacity=0.75]
  (t5.north) to[bend right=38] node[above,font=\scriptsize,pos=0.5]{0.55} (t3.north);
\draw[-{Stealth[length=2mm]},ourscol,line width=1.2pt,opacity=0.6]
  (t5.north) to[bend right=30] node[above,font=\scriptsize,pos=0.45]{0.25} (t4.north);
\draw[-{Stealth[length=2mm]},ourscol,line width=0.6pt,opacity=0.45]
  (t5.north) to[bend right=42] node[below,font=\scriptsize,pos=0.6]{0.12} (t2.north);
\draw[-{Stealth[length=2mm]},ourscol,line width=0.4pt,opacity=0.35]
  (t5.north) to[bend right=46] (t1.north);

\node[text width=13cm,align=center,font=\scriptsize,below=1.1cm of t3]
  {Each arc is an \emph{attention weight}: how much this position draws on that
   one. The weights are computed from the content, not fixed in advance, and
   they sum to 1. Here the model attends most strongly to the poly-alanine
   block --- plausibly because a $\beta$-sheet region several tokens back
   constrains what can sensibly follow.};
\end{tikzpicture}
\caption[Attention]{Self-attention at one position. Thickness indicates weight.
Because weights are recomputed from content at every position, the model can
route information over arbitrary distances --- which is what lets it connect a
crystalline block to a distant terminal domain.}
\label{fig:attention}
\end{figure}

Three further pieces complete the picture:

\begin{description}[leftmargin=0cm,style=nextline,itemsep=4pt]
\item[Multi-head attention.] The above is done several times in parallel with
  different learned projections --- SilkomeGPT uses \paperterm{8 heads}. One head
  might track repeat periodicity, another composition, another distance to the
  terminal domain. They are concatenated.

\item[Layers.] The whole operation (attention, then a small feed-forward
  network) is stacked. SilkomeGPT has \paperterm{12 layers}. Early layers
  handle local patterns; later layers combine them into longer-range structure.

\item[Positional information.] Attention alone is order-blind --- it sees a bag
  of tokens. Position must be injected. SilkomeGPT uses \term{rotary positional
  embeddings} (RoPE), which encode position by rotating each token's vector by
  an angle proportional to its index. The useful property is that the
  \emph{relative} offset between two positions is what survives into the
  attention computation, so patterns learned at one place transfer to another
  --- well suited to a protein made of the same motif repeated at many offsets.
\end{description}

\begin{figure}[htbp]
\centering
\begin{tikzpicture}[font=\small,
  lay/.style={draw=papercol!60,fill=papercol!8,rounded corners=2pt,
              minimum width=6.2cm,minimum height=0.62cm,inner sep=3pt},
  sub/.style={draw=greyc!60,fill=boxbg,rounded corners=2pt,
              minimum width=2.7cm,minimum height=0.55cm,font=\scriptsize},
  ar/.style={-{Stealth[length=2mm]},greyc!80},
]
\node[lay] (emb) at (0,0) {token embedding \;+\; rotary position};
\node[lay,above=0.45cm of emb] (l1) {transformer block 1};
\node[lay,above=0.28cm of l1] (l2) {transformer block 2};
\node[above=0.20cm of l2,font=\small] (dots) {$\vdots$};
\node[lay,above=0.20cm of dots] (l12) {transformer block 12};
\node[lay,above=0.45cm of l12,fill=ourscol!10,draw=ourscol!60] (head)
  {output layer $\rightarrow$ probability over 50{,}000 tokens};

\draw[ar] (emb) -- (l1);
\draw[ar] (l1) -- (l2);
\draw[ar] (l12) -- (head);

% expansion of one block
\node[sub] (att) at (6.4,1.15) {8-head self-attention};
\node[sub,below=0.35cm of att] (ff) {feed-forward, width 4096};
\draw[ar] (att) -- (ff);
\node[draw=greyc!45,dashed,rounded corners=3pt,
      fit=(att)(ff),inner sep=6pt] (blk) {};
\node[above=1pt of blk,font=\scriptsize,text=greyc] {inside one block};
\draw[greyc!60,dashed] (l2.east) -- (blk.west);

\node[below=0.45cm of emb,font=\scriptsize,text width=6.4cm,align=center]
  {input: \code{CalculateSilkContent<GAGAGS\dots>}};
\node[above=0.35cm of head,font=\scriptsize,text width=6.4cm,align=center]
  {output: the next token, sampled or taken greedily};
\end{tikzpicture}
\caption[Architecture]{SilkomeGPT's architecture. Hidden width 1024,
feed-forward width 4096, 12 blocks, 8 heads --- 253{,}556{,}736 parameters in
total. ``Decoder-only'' means there is just this one stack: no separate encoder,
and each position may attend only to positions before it.}
\label{fig:arch}
\end{figure}

\begin{keybox}{Parameters}
A \term{parameter} is one learned number. SilkomeGPT has 253.6 million of them.
Training means adjusting all of them so that the model's predicted
next-token probabilities match what actually came next in the training data.
Verifying that the released model really has this many, and in this
arrangement, was the first check I ran (Section~\ref{sec:checkpoint}).
\end{keybox}

\section{Pretraining and fine-tuning}

Training happened in two stages, and the distinction matters for interpreting
the results.

\begin{description}[leftmargin=0cm,style=nextline,itemsep=5pt]
\item[Pretraining.] The model saw $\approx$600{,}000 protein sequences from
  across biology and learned to predict the next token. No properties involved.
  The point is to absorb the general statistics of proteins --- which residues
  follow which, what real sequences look like.

\item[Fine-tuning.] The model was then trained further on the 1{,}033 silk
  sequences \emph{paired with their measured properties}, in a specific text
  format. This is where the property knowledge enters.
\end{description}

\begin{warnbox}{The fact that shapes every later interpretation}
The paper states that fine-tuning used ``all known pairs of sequence and
properties''. There is no \term{held-out set} --- no sequences kept aside to
test on. Every one of the 1{,}033 pairs was seen during training.

This is not a mistake; with 1{,}033 examples you want to use all of them, and
the paper reports self-consistency rather than predictive accuracy. But it means
no number measured on those 1{,}033 pairs can distinguish \emph{learning} from
\emph{memorising}. Chapter~\ref{ch:extensions} shows that the distinction is
not academic here.
\end{warnbox}

\section{Decoding: turning probabilities into text}
\label{sec:decoding}

The model outputs a probability for each possible next token. Turning that into
an actual choice is called \term{decoding}, and the settings involved caused
more subtlety in this project than anything in the architecture.

\begin{description}[leftmargin=0cm,style=nextline,itemsep=5pt]
\item[Greedy decoding.] Always take the highest-probability token.
  Deterministic: the same input always gives the same output.

\item[Sampling.] Draw randomly according to the probabilities. Different every
  time. Necessary for generating variety.

\item[Temperature $T$.] Rescales the distribution before sampling.
  $T \to 0$ approaches greedy; $T = 1$ samples from the raw distribution;
  $T > 1$ flattens it, making unlikely tokens more likely. The paper's inverse
  task uses $T = 1.25$ --- deliberately adventurous, because it wants novel
  designs.

\item[top-$k$ and top-$p$.] Truncation rules. top-$k$ keeps only the $k$ most
  likely tokens; top-$p$ (nucleus sampling) keeps the smallest set whose
  probabilities sum to $p$. Both stop the model from occasionally emitting
  something absurd.
\end{description}

\begin{figure}[htbp]
\centering
\begin{tikzpicture}[font=\small]
\begin{axis}[
  width=13cm,height=4.6cm,
  ybar, bar width=7pt,
  ymin=0,ymax=0.62, xmin=0.3, xmax=10.7,
  xtick={1,...,10},
  xticklabels={\aaseq{A},\aaseq{G},\aaseq{S},\aaseq{Q},\aaseq{Y},\aaseq{L},\aaseq{P},\aaseq{V},\aaseq{R},\aaseq{W}},
  ylabel={probability},
  axis lines=left, ylabel style={font=\scriptsize},
  legend style={font=\scriptsize,draw=none,at={(0.99,0.97)},anchor=north east},
]
\addplot[fill=papercol!30,draw=papercol] coordinates
  {(1,0.30)(2,0.22)(3,0.14)(4,0.11)(5,0.08)(6,0.06)(7,0.04)(8,0.03)(9,0.015)(10,0.005)};
\addlegendentry{$T=1$ (raw)}
\addplot[fill=ourscol!35,draw=ourscol] coordinates
  {(1,0.56)(2,0.25)(3,0.09)(4,0.05)(5,0.025)(6,0.012)(7,0.007)(8,0.004)(9,0.001)(10,0.0003)};
\addlegendentry{$T=0.5$ (sharper)}
\end{axis}
\end{tikzpicture}
\caption[Temperature]{Temperature reshapes the next-token distribution.
Lowering $T$ concentrates mass on the leader, making output more predictable
and less varied; raising it does the opposite. At $T \to 0$ sampling becomes
greedy decoding.}
\label{fig:temperature}
\end{figure}

\begin{keybox}{The settings used in this project}
\begin{center}
\small
\begin{tabular}{lcccc}
\toprule
& temperature & top-$k$ & top-$p$ & max new tokens \\
\midrule
inverse task (design) & 1.25 & 500 & 0.8 & 512 \\
forward task (predict), \paperterm{paper} & 0.01 & 500 & 0.9 & 64 \\
forward task (predict), \ourterm{ours} & \multicolumn{3}{c}{greedy} & 64 \\
\bottomrule
\end{tabular}
\end{center}
None of these appear in the paper; all were read out of the authors' released
notebook. Why I departed from the paper on the forward task, and the
measurements that justified it, is Section~\ref{sec:decodingchoice}.
\end{keybox}
